% Глава 3. Метод и экспериментальная установка
% Подраздел 3.1. Теоретические основы предлагаемого подхода

\subsection{Теоретические основы предлагаемого подхода}
\label{sec:method}

\subsubsection*{Архитектура FLUX.1}

FLUX.1~\cite{blackforest2024flux}~--- семейство диффузионных трансформеров Black~Forest~Labs, построенных на основе многомодального DiT (MM-DiT)~\cite{esser2024sd3,peebles2023dit}. Используется вариант FLUX.1-dev (\texttt{black-forest-labs/FLUX.1-dev}, точная ревизия зафиксирована в конфигурациях; вычисления в bfloat16 на NVIDIA~H100).

Конвейер включает: вариационный автоэнкодер~\cite{rombach2022ldm} со сжатием 8$\times$ в 16-канальное латентное пространство; два текстовых энкодера~--- CLIP~ViT-L/14~\cite{radford2021clip} (77~токенов, pooled projection как условный вектор) и T5-XXL~\cite{raffel2020t5} (до 256~токенов, участвует в joint attention); планировщик \texttt{FlowMatchEulerDiscrete}~\cite{von2022diffusers} при 28~шагах и CFG-guidance~$3{,}5$~\cite{ho2022cfg}. Ключевая особенность~--- разделение бэкбона: первые 19~блоков двухпоточные (double-stream, DS), визуальные и текстовые токены проходят независимыми проекциями с последующим joint attention; следующие 38~блоков однопоточные (single-stream, SS) с конкатенацией модальностей в единый self-attention.

\begin{figure}[h]
    \centering
    \includegraphics[width=0.95\textwidth]{flux_blora_architecture.png}
    \caption{Архитектура FLUX.1-dev и расположение блочных адаптеров B-LoRA. Стилевой адаптер $\Theta_\text{style}$ обучается на ранних double-stream блоках ($[0\text{--}18]$ в основной конфигурации; на схеме акцентированы ds\_00--ds\_01 с наивысшей style-специфичностью по данным Phase~0); содержательный адаптер $\Theta_\text{content}$~--- на поздних double-stream блоках $[9\text{--}18]$. Single-stream блоки в предлагаемом методе не модифицируются.}
    \label{fig:flux_architecture}
\end{figure}

В качестве генеративного транспорта используется метод спрямлённых потоков (rectified flow)~\cite{liu2022rectifiedflow}. В настоящей главе $x_0 \sim \mathcal{N}(0, \mathbf{I})$~--- шум, $x_1$~--- данные, $t \in [0, 1]$~--- время транспорта (нотация противоположна разделу~\ref{sec:literature_review}). Прямолинейный транспорт задаётся ОДУ:
\begin{equation}
    dx = (x_1 - x_0)\,dt,
    \label{eq:rectflow}
\end{equation}
а трансформер обучается предсказывать поле скоростей $(x_1 - x_0)$ минимизацией:
\begin{equation}
    \mathcal{L}(\theta) = \mathbb{E}_{t,\,x_0,\,x_1}\!\left[\bigl\|f_\theta\!\left((1-t)x_0 + t\,x_1,\;t\right) - (x_1 - x_0)\bigr\|^2\right].
    \label{eq:flow_loss}
\end{equation}
Формализм допускает меньшее число шагов интегрирования по сравнению с DDPM и упрощает масштабирование на высокие разрешения.

\subsubsection*{LoRA: адаптация низкого ранга}

LoRA~\cite{hu2022lora} вводит аддитивное обновление низкого ранга для замороженной матрицы весов $W_0 \in \mathbb{R}^{d \times k}$:
\begin{equation}
    W = W_0 + BA, \quad B \in \mathbb{R}^{d \times r},\; A \in \mathbb{R}^{r \times k},\; r \ll \min(d, k).
    \label{eq:lora}
\end{equation}
Полное число обучаемых параметров $r(d+k)$ вместо $dk$: при $r{=}16$ и типичных $d{=}k{=}3072$ для FLUX.1 экономия примерно в 96~раз. Адаптация применяется к проекциям $Q$, $K$, $V$ и \texttt{out\_proj} в каждом блоке. Инициализация $A \sim \mathcal{N}(0, \sigma^2)$, $B = 0$ гарантирует $\Delta W = 0$ в начале обучения. На инференсе вводится множитель силы адаптера $s$:
\begin{equation}
    W = W_0 + s \cdot \frac{\alpha_\text{train}}{r}\, BA,
    \label{eq:lora_scale}
\end{equation}
здесь $\alpha_\text{train}{=}r$ (единичная инициализация по \texttt{peft}); далее $\alpha$ обозначает именно $s$~--- множитель применения адаптера при инференсе.

\subsubsection*{B-LoRA в оригинальной постановке (SDXL)}

Метод B-LoRA~\cite{frenkel2024blora} опирается на эмпирическое наблюдение в SDXL: блок \texttt{up\_blocks[0].attentions[1]} (\textit{W2}, Up-декодер) преимущественно кодирует стиль~--- текстуру, цветовую палитру, художественную манеру,~--- тогда как \texttt{down\_blocks[2].attentions[1]} (\textit{W1}, Down-энкодер) отвечает за структурное содержание. Файнтюн на одном референсе с DreamBooth-промптом~\cite{ruiz2023dreambooth} вида <<\texttt{a [v] in [s] style}>> даёт два независимых адаптера $\Theta_\text{style}$ и $\Theta_\text{content}$, применимых отдельно или совместно. Разделение стиля и содержания достигается имплицитно за счёт структурной специализации блоков, без явных регуляризирующих потерь.

\subsubsection*{Анализ специализации блоков FLUX.1 (Phase~0)}

Поскольку MM-DiT не содержит U-Net-блоков, функциональные роли определяются эмпирически. Предлагается методология \textit{prompt injection}: для каждого блока $b$ текстовый эмбеддинг T5-XXL заменяется эмбеддингом альтернативного промпта, прочие входы остаются неизменными. Степень реакции сгенерированного изображения служит мерой того, насколько блок $b$ «прислушивается» к текстовому условию.

Эксперимент проводится в двух режимах: \textit{style injection} (базовый $P_\text{content}=\texttt{"A photo of a \{object\}"}$, инъецируемый $P_\text{style}=\texttt{"A photo of a \{color\} \{object\}"}$) и \textit{content injection} (направление обратное). $N{=}200$~пар промптов, все 19~DS и все 38~SS-блоки. Мерой реакции служит:
\begin{equation}
    S_b = \frac{1}{N} \sum_{i=1}^{N}
        \cos\!\bigl(\mathrm{CLIP}(I_b^{(i)}),\; \mathrm{CLIP}_\text{text}(P_\text{inject}^{(i)})\bigr),
    \label{eq:phase0_score}
\end{equation}
где $I_b^{(i)}$~--- изображение, сгенерированное с инъекцией в блок~$b$. Высокое $S_b$ указывает на активное участие блока в трансляции текстового условия.

\begin{table}[h]
    \centering
    \small
    \caption{Чувствительность выбранных блоков FLUX.1-dev к инъекции промпта ($N=200$, CLIP ViT-L/14). Специфичность $= S_\text{style} - S_\text{content}$.}
    \label{tab:phase0_results}
    \begin{tabular}{llrrr}
        \toprule
        Блок & Тип & $S_\text{style}$ & $S_\text{content}$ & Специфичность \\
        \midrule
        ds\_00 & DS & 0.260 & 0.243 & \textbf{+0.018} \\
        ds\_01 & DS & 0.229 & 0.225 & +0.004 \\
        ds\_02 & DS & 0.187 & 0.189 & $-$0.003 \\
        ds\_05 & DS & 0.129 & 0.146 & $-$0.017 \\
        ds\_12 & DS & 0.133 & 0.154 & $-$0.021 \\
        ds\_18 & DS & 0.163 & 0.159 & +0.003 \\
        \midrule
        ss\_06 & SS & 0.175 & 0.176 & $-$0.001 \\
        ss\_07 & SS & 0.221 & 0.226 & $-$0.004 \\
        ss\_10 & SS & 0.226 & 0.220 & +0.005 \\
        ss\_23 & SS & 0.242 & 0.240 & +0.002 \\
        ss\_28 & SS & 0.251 & 0.254 & $-$0.002 \\
        ss\_37 & SS & 0.230 & 0.236 & $-$0.006 \\
        \bottomrule
    \end{tabular}
\end{table}

\begin{figure}[h]
    \centering
    \includegraphics[width=\textwidth]{block_analysis_chart.png}
    \caption{Полная карта чувствительности всех 19~double-stream и 38~single-stream блоков FLUX.1-dev по результатам Phase~0 (CLIP-embedding injection, $N=200$~пар промптов). Видны три закономерности: фазовая граница ds\_04$\to$ds\_05, плато single-stream блоков, начиная с ss\_07, и локальные провалы ss\_20/ss\_29.}
    \label{fig:phase0_block_chart}
\end{figure}

\textbf{Граница DS~early\,/\,DS~mid.} Блоки ds\_00 и ds\_01 образуют зону наибольшей style-специфичности; на переходе ds\_04$\to$ds\_05 фиксируется резкий спад $S_\text{style}$ ($\Delta\approx 0.057$). Блоки ds\_05--ds\_12 образуют content-доминирующую долину с экстремумами $-0.021$ (ds\_12) и $-0.017$ (ds\_05).

\textbf{Граница SS~early\,/\,SS~plateau.} Блоки ss\_00--ss\_06 сопоставимы с DS~mid ($S_b\approx 0.17\text{--}0.20$); на переходе ss\_06$\to$ss\_07 происходит скачок $\Delta\approx 0.047$, после чего 31~оставшийся блок образует однородное плато ($S_b\approx 0.22\text{--}0.25$). Style/content-специфичность SS-блоков в плато статистически незначима.

\textbf{Аномальные блоки.} ss\_20 и ss\_29 демонстрируют локальные провалы ($S_b$ ниже соседей на 0.022--0.034) одновременно в обоих режимах, что свидетельствует об архитектурной особенности FLUX.1 и исключает эти блоки из аблационных диапазонов.

\textbf{Выбор блоков.} Кандидаты в стилевые~--- DS~$[0\text{--}18]$ (полный диапазон тестируется в аблации D-A); кандидаты в содержательные (Phase~1b)~--- три перекрывающихся подмножества double-stream: ранний $[0\text{--}8]$ (DC-01), поздний $[9\text{--}18]$ (DC-02) и центральная зона $[5\text{--}12]$ (DC-03). Итоговый выбор $\Theta_\text{content}$ определяется по метрикам обучения.

\subsubsection*{Адаптация B-LoRA к FLUX.1}

Прямой перенос B-LoRA на FLUX.1 невозможен: модель не содержит U-Net и блоков \texttt{up\_blocks}/\texttt{down\_blocks}. Функциональные аналоги стилевых и содержательных блоков ищутся в пространстве MM-DiT экспериментально, а не по архитектурной аналогии.

\textbf{Рабочая гипотеза:} double-stream блоки $[0\text{--}18]$ FLUX.1, реализующие joint attention между текстовым и визуальным потоками, функционально соответствуют стилевому блоку \textit{W2} Up-декодера SDXL~\cite{frenkel2024blora}. Гипотеза согласуется с работами FluxSpace~\cite{dalva2024fluxspace} и SplitFlux~\cite{yang2025splitflux}, показавших функциональную специализацию блоков MM-DiT.

По результатам Phase~1 и Phase~1b в контрольной точке GATE~1 принимается один из трёх сценариев:
\begin{itemize}\itemsep0.2em
    \item \textbf{Сценарий~A}~--- DS-блоки дают перенос стиля \textit{и} выделен $\Theta_\text{content}$: полное разделение стиля и содержания (дополнительно выполняется Phase~6 style-content mixing);
    \item \textbf{Сценарий~B}~--- DS-блоки дают перенос стиля, но $\Theta_\text{content}$ не выделяется: применяется только $\Theta_\text{style}$, Phase~6 пропускается;
    \item \textbf{Сценарий~C}~--- DS-блоки не обеспечивают переноса: проводится дополнительный анализ альтернативных подмножеств, конфигурация переопределяется.
\end{itemize}

\begin{figure}[h]
    \centering
    \includegraphics[width=0.98\textwidth]{gate1_decision.png}
    \caption{Структура решения в контрольной точке GATE~1.
             Phase~1 (серия~D на DS~$[0\text{--}18]$) и Phase~1b (серия~DCC,
             поиск $\Theta_\text{content}$) образуют двумерный критерий:
             значимость стилевого сдвига и наличие выделяемого содержательного
             адаптера. По итогам работы (раздел~\ref{sec:results_phase1b})
             принят Сценарий~A: $\Theta_\text{style}{=}$DS~$[0\text{--}18]$,
             $\Theta_\text{content}{=}$DS~$[9\text{--}18]$.}
    \label{fig:gate1_decision}
\end{figure}

\textbf{Реализация метода.} (1)~Все веса FLUX.1 замораживаются. (2)~К double-stream блокам $[0\text{--}18]$ добавляется LoRA-адаптер ранга~$r$~--- $\Theta_\text{style}$; выделение $\Theta_\text{content}$ производится при Сценарии~A. (3)~Адаптер обучается на одном референсе за $T$~шагов (из $\{500,\,1\,000,\,2\,000,\,4\,000\}$, оптимум по~\ref{sec:ablations}) с целевой функцией flow-matching~\eqref{eq:flow_loss}, скоростью обучения $5{\times}10^{-5}$, оптимизатором AdamW~\cite{loshchilov2019adamw} ($\beta_1{=}0{,}9$, $\beta_2{=}0{,}999$, weight decay $10^{-2}$, без warmup, без EMA), в bfloat16 средствами Hugging~Face \texttt{accelerate}~\cite{accelerate2022} и фреймворка \texttt{ai-toolkit}. (4)~При инференсе подключается $\Theta_\text{style}$: модель воспроизводит стиль референса, а содержание задаётся текстовым промптом.

\subsubsection*{Поиск адаптера содержания $\Theta_\text{content}$ (Phase~1b)}

В оригинальном B-LoRA~\cite{frenkel2024blora} $\Theta_\text{style}$ и $\Theta_\text{content}$ приписаны заранее известным блокам SDXL; в FLUX.1 такого априорного соответствия нет. По данным Phase~0 наибольшая content-специфичность сосредоточена в диапазоне ds\_05--ds\_12. В Phase~1b обучаются и сравниваются три подмножества: ранний $[0\text{--}8]$ (DC-01), поздний $[9\text{--}18]$ (DC-02), центральный $[5\text{--}12]$ (DC-03). На каждом диапазоне обучается отдельный LoRA-адаптер ($r{=}16$, прочие гиперпараметры как в Phase~1); при инференсе применяется только этот адаптер без $\Theta_\text{style}$.

Адаптер признаётся содержательным, если относительно базовой модели достигается (i)~рост CLIP-content или (ii)~статистически значимое отклонение DINO-style с ограничением подъёма фактором~$<\!5$. При выполнении условия хотя бы для одного диапазона фиксируется \textbf{Сценарий~A}, иначе~--- \textbf{Сценарий~B}; обработка \textbf{Сценария~C} описана выше.
